\label{sec:marco-teorico}
\section{Ecuaciones diferenciales como problemas de aproximación}
El objeto de estudio es cómo distribuir recursos de representación para aproximar soluciones de ecuaciones diferenciales parciales (PDEs, por sus siglas en inglés). Una PDE relaciona magnitudes físicas y derivadas; la solución depende conjuntamente del operador, dominio, parámetros y condiciones auxiliares. No se presupone que toda PDE carezca de solución cerrada ni que un aproximador neuronal sea preferible a un solucionador clásico \cite{evans1998pde,raissi2019,grossmann2024}.

Sea $\Omega\subset\mathbb R^d$ el dominio espacial, $\partial\Omega$ su frontera, $t\in[0,T]$ el tiempo y $\mathbf u:Q\to\mathbb R^m$ el conjunto de campos. Para problemas estacionarios $Q=\Omega$; para los evolutivos $Q=\Omega\times(0,T]$. Escribiremos $z=x$ o $z=(x,t)$, respectivamente:
\begin{equation}
\mathcal N_a[\mathbf u;\mu](z)=f_a(z),\qquad \mathcal B_b[\mathbf u](z)=g_b(z),\qquad \mathbf u(x,0)=\mathbf u_0(x).
\label{eq:pde-general}
\end{equation}
$a$ identifica ecuaciones, $b$ condiciones y $\mu$ parámetros físicos. $\mathcal B_b$ impone valores (Dirichlet), flujos (Neumann) o combinaciones (Robin). La condición inicial se omite en problemas estacionarios. Una ecuación puede involucrar varios campos: índice de ecuación e índice de campo no son intercambiables.

La aproximación $\mathbf u_\theta$ pertenece a una familia determinada por arquitectura y parámetros $\theta$. Se distinguen error $e=\mathbf u_\theta-\mathbf u$ y residual
\begin{equation}
r_a(z;\theta)=\mathcal N_a[\mathbf u_\theta;\mu](z)-f_a(z).
\label{eq:residual}
\end{equation}
El forzamiento se mantiene explícito. Un residual pequeño en puntos finitos no certifica una solución precisa: intervienen estabilidad, frontera, cobertura y optimización. Usaremos $\|v\|_{L^2(Q)}^2=\int_Q|v|^2dz$ y $|v|_{H^1(\Omega)}^2=\int_\Omega\|\nabla v\|^2dx$; la segunda detecta discrepancias de gradiente.

Las soluciones analíticas permiten verificar operadores sin incertidumbre de discretización. FDM, FVM y FEM proporcionan referencias cuando no hay solución cerrada. El aprendizaje automático científico (SciML) integra física y aprendizaje; su utilidad depende de la tarea y del costo total, incluido el entrenamiento \cite{karniadakis2021,grossmann2024}. La referencia independiente fija qué significa resolver correctamente el problema.

\section{Aproximación clásica y resolución adaptativa}
Una discretización convierte funciones en incógnitas finitas. La malla describe puntos, elementos y relaciones geométricas; tamaño y orden de aproximación controlan las escalas representables. Es el precedente de la adaptación neuronal, sin transferir automáticamente sus garantías.

\subsection{Diferencias y volúmenes finitos}
En diferencias finitas (FDM), Taylor da $u(x_i\pm h)=u_i\pm hu_i'+h^2u_i''/2\pm h^3u_i'''/6+O(h^4)$. Sumando y despejando,
\begin{equation}u''(x_i)=\frac{u_{i-1}-2u_i+u_{i+1}}{h^2}+O(h^2).\end{equation}
La precisión requiere suavidad y resolución: iterar más no resuelve una capa que la discretización no representa. En volúmenes finitos (FVM), integrar $u_t+\nabla\cdot F(u)=s$ sobre una celda $K$ y aplicar divergencia produce
\begin{equation}\frac d{dt}\int_Ku\,dx+\sum_{e\subset\partial K}\int_e\widehat F\cdot n_K\,ds=\int_Ks\,dx.\end{equation}
El flujo numérico $\widehat F$ compartido entre celdas permite balances discretos conservativos. FDM aproxima derivadas; FVM aproxima balances. Ambos admiten adaptación.

\subsection{Construcción de elementos finitos}
El método de elementos finitos (FEM, por sus siglas en inglés) aproxima con funciones asociadas a elementos. Considérese $-\nabla\cdot(\kappa\nabla u)=f$, $\kappa>0$, $u=0$ en $\Gamma_D$, $\kappa\nabla u\cdot n=g_N$ en $\Gamma_N$. Multiplicar por una prueba $v$ nula en $\Gamma_D$, integrar y aplicar integración por partes da
\begin{equation}
\underbrace{\int_\Omega\kappa\nabla u\cdot\nabla v\,dx}_{a(u,v)}=\underbrace{\int_\Omega fv\,dx+\int_{\Gamma_N}g_Nv\,ds}_{\ell(v)}.
\label{eq:weak}
\end{equation}
Se busca $u\in V=\{v\in H^1(\Omega):v|_{\Gamma_D}=0\}$. Con Dirichlet no homogéneo se usa un levantamiento y un espacio afín. Una malla $\mathcal T_h$ define $V_h\subset V$. Sustituir $u_h=\sum_jU_j\phi_j$ y probar con cada $\phi_i$ produce
\begin{equation}KU=F,\quad K_{ij}=\int_\Omega\kappa\nabla\phi_j\cdot\nabla\phi_i\,dx,\quad F_i=\ell(\phi_i).\end{equation}
En un triángulo lineal, $\phi_j$ es afín, vale uno en su nodo y cero en los restantes. Su soporte físico comprende los elementos que comparten el nodo. Esta construcción explica la asignación local de grados de libertad \cite{fenics2016}.
\begin{figure}[htbp]\centering\includegraphics[width=\linewidth]{figuras/fem_poisson.png}
\caption{Cálculo ilustrativo propio con FEM triangular $P_1$: malla, solución y error absoluto para $u=\sin(\pi x)\sin(\pi y)$. Se ensambla y resuelve el sistema; no es un resultado del mecanismo ni una simulación piezoeléctrica.}\label{fig:fem}\end{figure}

\subsection{Estimación a posteriori y refinamiento hp}
Un estimador a posteriori utiliza la solución discreta. Para Poisson de coeficiente unitario con elementos conformes, un indicador ilustrativo es
\begin{equation}
\eta_K^2=h_K^2\|f+\Delta u_h\|_{L^2(K)}^2+\tfrac12\sum_{e\subset\partial K\cap\Omega}h_e\|[\nabla u_h\cdot n_e]\|_{L^2(e)}^2.
\label{eq:estimator}\end{equation}
$h_K$ es diámetro, $h_e$ longitud de arista y los corchetes denotan salto con orientación común. Con Neumann se añaden discrepancias de frontera. Las cotas de error requieren hipótesis sobre malla, operador y datos; el indicador no es una identidad con el error. La marcación de Dörfler selecciona $\mathcal M$ con $\sum_{K\in\mathcal M}\eta_K^2\geq\vartheta\sum_K\eta_K^2$, $0<\vartheta<1$ \cite{dorfler1996}. Refinar $h$ subdivide elementos; $p$ aumenta grado polinómico; $hp$ combina ambos según regularidad. El ciclo es resolver, estimar, marcar y adaptar. La malla puede asignar resolución local, pero redes con descomposición o bases adaptativas también disponen de instrumentos de localidad.

\begin{figure}[htbp]\centering\includegraphics[width=.68\linewidth]{figuras/fem_convergencia.png}\caption{Convergencia comprobada de la demostración FEM en tres mallas. No valida el controlador futuro.}\label{fig:convergencia}\end{figure}

\section{Aprendizaje informado por física y estado del arte}
El estado del arte (SOTA, por sus siglas en inglés) se organiza por el objeto modificado: evaluación, pérdida, representación o entrenamiento. Esta sección establece conceptos y comparadores; la pregunta posterior aprovecha el mapa sin repetir el catálogo.

\subsection{De inteligencia artificial a ANN, DNN y MLP}
Inteligencia artificial engloba sistemas de percepción, inferencia o decisión. El aprendizaje automático (ML) ajusta modelos a ejemplos u objetivos; el aprendizaje profundo (DL) usa composiciones de varias capas. Una red neuronal artificial (ANN) es una familia de aproximadores; una red profunda (DNN) utiliza varias capas transformadoras. Ninguna denominación implica por sí sola información física.

El perceptrón multicapa (MLP) conecta capas de neuronas. Cada neurona combina una transformación afín y una activación no lineal. Con $h^{(0)}=z$,
\begin{equation}a_j^{(\ell+1)}=\sum_iW_{ji}^{(\ell)}h_i^{(\ell)}+b_j^{(\ell)},\qquad h_j^{(\ell+1)}=\sigma(a_j^{(\ell+1)}).\end{equation}
$W,b$ se aprenden; la forma de $\sigma$ suele fijarse, por ejemplo $\tanh$. Para residuales fuertes de segundo orden se necesitan derivadas apropiadas: una ReLU ordinaria tiene segunda derivada nula casi en todas partes. La retropropagación aplica la regla de la cadena a parámetros. La diferenciación automática deriva la función implementada; no elimina errores de modelo o aproximación \cite{lagaris1998,raissi2019}.
\begin{figure}[htbp]\centering\includegraphics[width=\linewidth]{figuras/arquitecturas.png}\caption{MLP con pesos en aristas y activación en nodos; KAN con funciones aprendibles en aristas. Ambos pueden utilizar una pérdida física.}\label{fig:arquitecturas}\end{figure}

\subsection{Redes neuronales informadas por física (PINNs)}
Una PINN aproxima $\mathbf u$ e incorpora ecuaciones, condiciones y, cuando existen, observaciones. Sean $R_a>0$ escalas físicas de ecuaciones y $U_v>0$ escalas de campos. Se construyen $\widetilde r_a=r_a/R_a$ y discrepancias adimensionales de frontera $\widetilde b_b$ e iniciales $\widetilde i_v$:
\begin{equation}
\mathcal L=\sum_a\frac{\lambda_a}{N_r}\sum_i\widetilde r_a(z_i)^2+\sum_b\frac{\lambda_b^B}{N_b}\sum_j\widetilde b_b(z_j)^2+\sum_v\frac{\lambda_v^I}{N_0}\sum_k\widetilde i_v(x_k)^2.
\label{eq:perdida}\end{equation}
Los puntos de colocación son posiciones de evaluación del operador, no elementos FEM. Las escalas se fijan antes de comparar métodos; normalizar por el error actual podría ocultar degradación. Datos supervisados añadirían un término declarado. Las PINNs admiten problemas directos e inversos, pero su precisión y costo dependen de condicionamiento, arquitectura y entrenamiento \cite{raissi2019,wang2023expertguide}.
\begin{figure}[htbp]\centering\includegraphics[width=\linewidth]{figuras/perdida_pinn.png}\caption{Construcción de la pérdida física. La referencia para evaluar error permanece fuera del controlador y del entrenamiento.}\label{fig:loss}\end{figure}
El sesgo espectral describe tendencias de aprendizaje por frecuencia bajo ciertas redes y condiciones \cite{rahaman2019}; no demuestra incapacidad de representar altas frecuencias ni se explica únicamente por ausencia de malla. Se separarán errores de representación, optimización y discretización de la pérdida.

\subsection{KAN y PIKAN: construcción de B-splines}
El teorema de representación de Kolmogorov--Arnold establece que una función continua $f:[0,1]^d\to\mathbb R$ admite
\begin{equation}f(x)=\sum_{q=1}^{2d+1}\Phi_q\!\left(\sum_{p=1}^d\psi_{qp}(x_p)\right)\end{equation}
con funciones univariadas continuas apropiadas. Motiva una estructura, pero no garantiza que una KAN finita, suave y entrenada numéricamente encuentre la representación. Una capa práctica es
\begin{equation}h_j^{(\ell+1)}=\sum_i\varphi_{ji}^{(\ell)}(h_i^{(\ell)}),\quad\varphi(a)=\alpha b(a)+\sum_{k=0}^{n-1}c_kB_{k,p}(a).\label{eq:kan}\end{equation}
$b$ puede ser SiLU; $\alpha,c_k$ son parámetros y $p$ es grado. Una red de Kolmogorov--Arnold informada por física (PIKAN) utiliza una KAN en la pérdida física; existen variantes con otras bases \cite{liu2024kan,shukla2024fair}.

Dados nudos no decrecientes $t_0,\ldots,t_{n+p}$, Cox--de Boor construye
\begin{align}
B_{k,0}(a)&=\mathbf1_{[t_k,t_{k+1})}(a),\\
B_{k,p}(a)&=\frac{a-t_k}{t_{k+p}-t_k}B_{k,p-1}(a)+\frac{t_{k+p+1}-a}{t_{k+p+1}-t_{k+1}}B_{k+1,p-1}(a).\label{eq:spline}
\end{align}
Los términos con denominador cero valen cero; el extremo final se evalúa por continuidad. El soporte es $[t_k,t_{k+p+1}]$. En nudos simples la continuidad es $C^{p-1}$; multiplicidades la reducen. Derivar da
\begin{equation}B'_{k,p}=\frac p{t_{k+p}-t_k}B_{k,p-1}-\frac p{t_{k+p+1}-t_{k+1}}B_{k+1,p-1}.\end{equation}
Reaplicar la identidad construye segundas derivadas. Nudos delimitan intervalos, coeficientes ponderan bases y centros parametrizan otras familias: no son sinónimos. Insertar nudos puede enriquecer el espacio; moverlos puede cambiarlo sin aumentar dimensión.
\begin{figure}[htbp]\centering\includegraphics[width=\linewidth]{figuras/bases.png}\caption{B-splines cúbicas, una función construida con ellas y una base compacta trasladable. Ilustraciones matemáticas sin entrenamiento.}\label{fig:bases}\end{figure}
Liu et al. ya incluyen extensión de grilla y transferencia entre representaciones \cite{liu2024kan}. Actualizar posiciones por activaciones y aumentar intervalos son operaciones distintas. La mera existencia de una grilla adaptable no es una contribución nueva.

\subsection{Qué redistribuyen las estrategias existentes}
RAR y RAD modifican colocación mediante residual \cite{wu2023sampling}. Una densidad tipo RAD es $q\propto |r|^k/\mathbb E|r|^k+c$. Si $q$ es densidad respecto a una medida uniforme $\mu$, $\mathbb E_q[r^2/q]=\mathbb E_\mu[r^2]$; omitir la corrección cambia el funcional esperado. Muestreo y ponderación pueden estar relacionados sin ser implementaciones idénticas.

Para explicitar otras construcciones del SOTA, un residual variacional local es $R_{K,j}=\int_K r(z;\theta)v_j(z)\,dz$; integrar por partes según el operador reduce el orden de derivación del aproximador. Una pérdida hp-VPINN suma $|R_{K,j}|^2$ sobre elementos y pruebas, con condiciones de frontera declaradas. En entrenamiento causal, una forma habitual de peso es $w_n=\exp(-\epsilon_c\sum_{j<n}\mathcal L_j)$, donde $\mathcal L_j$ mide residual de tiempos anteriores: el peso responde al avance temporal, no modifica por sí mismo el espacio representacional. En el entrenamiento básico, $\theta_{k+1}=\theta_k-\eta\nabla_\theta\mathcal L$ ilustra el papel del gradiente; Adam añade estimaciones de momentos. Cambiar pesos, geometría o estado del optimizador son intervenciones distintas.

hp-VPINN incorpora pruebas locales y partición \cite{kharazmi2021hpvpinns}; entrenamiento causal controla contribuciones temporales \cite{wang2024causal}; FBPINN y FBKAN hacen explícita la descomposición de representación \cite{moseley2023,howard2024}. AB-PINNs también adapta y enriquece subdominios durante el entrenamiento \cite{abpinn2025}. Por tanto, no es defendible afirmar que ninguna red informada por física asigna capacidad local.
\begin{table}[htbp]\centering\small\caption{Objeto intervenido y control relevante; no es un ranking universal de precisión.}\label{tab:rejilla}
\begin{tabularx}{\linewidth}{p{.27\linewidth}XX}\toprule
Familia & Intervención & Control\\\midrule
RAR/RAD & Colocación por residual & Mismo aproximador y costo\\
hp-VPINN & Pruebas y partición & Formulación variacional\\
Causal & Pesos por tiempo & Ventanas comunes\\
FBPINN/FBKAN & Representación local & Descomposición explícita\\
KAN con extensión & Resolución spline & Igual crecimiento y transferencia\\
PIKAN adaptativa & Grilla y entrenamiento & Ablaciones por componente\\
Densidad/curvatura & Distribución de nudos & Señales comparables\\
Nudos libres/multinivel & Geometría y optimización & Conservación y costo\\\bottomrule
\end{tabularx}\end{table}
Rigas et al. estudian adaptación PIKAN y transición del entrenamiento \cite{rigas2024}; el framework de 2026 dirige nudos con densidades de importancia, incluida curvatura \cite{rigas2026}. Nudos libres y entrenamiento multinivel aportan precedentes sobre geometría y optimización \cite{actor2025,actor2026}. RA-PIKAN combina densificación residual y adaptación del entrenamiento \cite{leong2026}. Se usan como respaldo para estudiar rutas aledañas, no como razón para descartar adaptación local. Sus resultados publicados no equivalen a reproducciones realizadas en esta tesis.

\section{Firma matemática de la dificultad}
La dificultad representacional es relativa al aproximador, presupuesto y norma. Aquí se definen propiedades de soluciones, sin confundirlas con señales que el modelo necesariamente observa. Se emplearán campos y coordenadas adimensionales.
\subsection{Oscilación, gradiente y curvatura}
En un dominio periódico, $E_{>k_0}=\sum_{|k|>k_0}|\widehat u_k|^2$ describe energía en frecuencias altas; la localización requiere ventanas. Para una región $A$,
\begin{equation}G(A)=\int_A\|\nabla u\|^2dx,\qquad H(A)=\int_A\|D^2u\|_F^2dx.\end{equation}
Miden variación y curvatura, no error. Una cota local $\|u-I_hu\|_{L^2(K)}\leq Ch_K^{p+1}|u|_{H^{p+1}(K)}$, bajo regularidad y condiciones del interpolante, explica reducir resolución donde derivadas altas son grandes. No convierte la segunda derivada en estimador universal para splines cúbicas.

Para densidad $d\geq0$, $C(A)=\int_A d/\int_\Omega d$ se compara con $|A|/|\Omega|$. Denominador cero significa ausencia de señal. Firmas de la referencia sirven para evaluar; firmas de $u_\theta$ son observables que pueden omitir la estructura todavía no aprendida.
\subsection{Movimiento, no linealidad y acoplamiento}
La linealización $\mathcal N[u+\delta u]\approx\mathcal N[u]+D\mathcal N[u]\delta u$ muestra sensibilidad a perturbaciones. En sistemas, $D_{u_v}\mathcal N_a$ describe dependencias entre campos. Una interfaz regular $u(x,t)=\gamma$ tiene velocidad normal $V_n=-u_t/\|\nabla u\|$: su movimiento debe comprobarse, no inferirse del nombre de la PDE. En interfaces materiales los coeficientes pueden saltar y requieren transmisión por subdominios.

La cadena es propiedad física, firma matemática, exigencia representacional y observable. Curvatura y residual pueden compartir derivadas sin ser iguales ni independientes. Se medirá su información incremental, sin suponerla por notación.

\section{Pregunta y contribución posible}
La pregunta propuesta es: \emph{¿puede un controlador decidir dónde reasignar un presupuesto fijo de representación, usando utilidad marginal estimada y restricciones de conservación, para reducir error localizado sin trasladarlo a otras regiones o campos?}

El aporte candidato es la regla de decisión y su caracterización causal, apoyada en grillas y bases publicadas. Se estudiarán información adicional frente a residual/curvatura, beneficio de retirar o desplazar resolución, preservación de derivadas y reparto entre campos. Combinar componentes no asegura una contribución importante: se requiere beneficio reproducible frente a controles cercanos y delimitación de fallas.

La tesis conserva identidad de mecanismo; piezoelectricidad prueba acoplamiento. Permanecen tres rutas: ponderación para esfuerzo, nudos para resolución y centros compactos móviles como alternativa geométrica. La primera será control; las otras son candidatas. La elección final depende del piloto, no de una afirmación anticipada de prioridad. La búsqueda tiene corte el 6 de septiembre de 2026 y debe actualizarse antes de cerrar el manuscrito.

\section{Construcción matemática previa al mecanismo}
Se especifican objetos, invariantes y una regla candidata, todavía por implementar y falsar. El ciclo es estado, señal, asignación, transferencia, comprobación y aceptación o reversión.
\subsection{Esfuerzo, representación y localidad efectiva}
Con geometría $\xi$ y coeficientes $\theta$, el aproximador pertenece a $\mathcal U_\xi$. Ponderar $\sum_iw_ir_i^2$ conserva ese espacio y modifica optimización; cambiar $\xi$ modifica la familia disponible. Actualizar coeficientes locales no crea grados de libertad. Mover bases con cantidad fija redistribuye resolución sin aumentar dimensión total.

Para una spline de una arista interna, la regla de la cadena da
\begin{equation}\frac{\partial u_v(z)}{\partial c_k^{(\ell)}}=\frac{\partial u_v(z)}{\partial h_j^{(\ell+1)}}B_{k,p}(h_i^{(\ell)}(z)).\label{eq:localidad}\end{equation}
El soporte está en la coordenada latente; su preimagen física puede ser extensa o disconexa. Una base de primera capa puede cubrir una franja. La localidad física se medirá con sensibilidades. Para un bloque $\theta_j$,
\begin{equation}I_{jv}(z)=\|\partial_{\theta_j}(u_v/U_v)\|^2,\quad A_{jv}(R)=\frac{\int_R I_{jv}d\mu}{\int_QI_{jv}d\mu+\epsilon}.\end{equation}
$A$ diagnostica atribución en una parametrización fijada; no es invariante bajo reescalamientos arbitrarios.

\subsection{Nudos y centros móviles}
Para $G$ intervalos en $[a,b]$, posiciones ordenadas se parametrizan con
\begin{equation}\Delta_j=\delta+(b-a-G\delta)\frac{e^{\eta_j}}{\sum_ke^{\eta_k}},\quad t_j=a+\sum_{k=1}^j\Delta_k,\quad0<\delta<(b-a)/G.\end{equation}
Fija extremos y evita intervalos degenerados, sin garantizar buen condicionamiento. Una alternativa compacta es $\varphi(a)=\alpha b(a)+\sum_jc_j\rho((a-\xi_j)/h_j)$, donde
\begin{equation}\rho(q)=\begin{cases}\exp[-1/(1-q^2)],&|q|<1,\\0,&|q|\geq1.\end{cases}\end{equation}
Es suave y compacta; su soporte continúa en la coordenada de la arista. Centros libres por gradiente constituyen un comparador obligatorio para aislar el controlador.

\subsection{Señal de utilidad marginal}
Un residual alto puede reflejar optimización deficiente o representación insuficiente. Se propone una prueba corta: desde un mismo estado se crean una rama sin cambio y otra con acción $a_j$; ambas reciben $K_p$ pasos con iguales datos y presupuesto. En un conjunto diagnóstico independiente $D$ se evalúan ecuaciones y condiciones:
\begin{equation}
s_j=\frac{\mathcal L_D(\operatorname{Train}_{K_p}(\theta,\xi))-\mathcal L_D(\operatorname{Train}_{K_p}(\operatorname{Transfer}(\theta,\xi,a_j)))}{C_j+\epsilon}.
\label{eq:signal}\end{equation}
$C_j$ incluye ensayo y transferencia. $s_j$ estima beneficio a corto plazo por costo; no es error verdadero ni una garantía causal. Puede ser ruidoso y caro. $D$ no contiene soluciones de referencia; prueba final nunca entra aquí. Se compara contra residual, curvatura, constante y señal permutada. Una señal cara sin beneficio neto será un resultado negativo, no una razón para cambiar retrospectivamente los criterios.

\subsection{Asignación, tipos y presupuesto}
El estado $S=(\theta,\xi,\mathcal O,D)$ incluye optimizador $\mathcal O$. El extractor $E:S\mapsto\mathbb R^{J\times m}$ produce utilidades de acciones; $T:\mathbb R^{J\times m}\to\Delta^{Jm-1}$ produce una distribución y el asignador $A:(\Delta,S,B)\mapsto\xi'$ devuelve geometría factible. Para unidades $j$ y campos $v$,
\begin{equation}\pi_{jv}=(1-\alpha)\frac{e^{\widetilde s_{jv}/\tau}}{\sum_{k,w}e^{\widetilde s_{kw}/\tau}}+\frac\alpha{Jm}.\label{eq:density}\end{equation}
$\tau>0$ controla concentración y $0<\alpha\leq1$ exploración. Escalas de $\widetilde s$ se fijan en piloto. Señal constante produce distribución uniforme, no necesariamente la grilla histórica. La geometría se expresa sobre aristas; las regiones físicas son diagnóstico.

Con $B$ bases y mínimos $B_v^{\min}$ se asignan $B_v=B_v^{\min}+\operatorname{LR}_v((B-\sum_wB_w^{\min})\sum_j\pi_{jv})$, donde mayores restos $\operatorname{LR}$ conserva la suma entera. Esto exige cabezas o subredes identificables por campo. En una red compartida no se atribuye presupuesto a campos sin una definición adicional. La alternativa continua mantiene $B_v$ y desplaza nudos. Una regla general es
\begin{equation}\max_{\mathcal A}\sum_{a\in\mathcal A}\widehat{\Delta\mathcal L}_a-\lambda_{\rm mov}D(\xi',\xi),\quad P(\xi')\leq P_{\max},\quad C(\mathcal A)\leq C_{\max},\quad\varepsilon_{\rm tr}\leq\varepsilon_{\max}.\end{equation}
$P$ cuenta parámetros, $C$ costo y $D$ magnitud del cambio. Las ganancias no necesariamente son aditivas: inicialmente se acepta una acción por evento. Añadir una base requiere retirar otra o compararse con igual crecimiento.

\subsection{Transferencia y conservación}
Mover o retirar bases puede alterar la función aun sin entrenamiento. Se proyecta sobre la nueva representación:
\begin{equation}
\theta'=\arg\min_{\vartheta}\sum_{z\in P,v}\left[\left|\frac{u_{\vartheta,\xi',v}-u_{\theta,\xi,v}}{U_v}\right|^2+\beta\left\|\nabla_{\widetilde x}\frac{u_{\vartheta,\xi',v}-u_{\theta,\xi,v}}{U_v}\right\|^2\right].\label{eq:transfer}
\end{equation}
$P$ incluye baja señal para prevenir olvido. Ajustar coeficientes de una arista con entradas congeladas es lineal; una proyección global profunda puede no serlo. Se comprueban valores, derivadas requeridas, residual y frontera en puntos ajenos al ajuste. Si falla se revierte. Inserción en espacios anidados puede admitir transferencia exacta; movimiento arbitrario no. En el análisis primario se reinicia el optimizador completo en cada evento, con reinicios simulados en controles en el mismo calendario; transferir momentos será una extensión separada.

\subsection{Reducciones, costo y límites}
Sin adaptación se recuperan arquitectura inicial y entrenamiento. Residual con cuantiles corresponde a una ruta próxima a adaptación por densidad. Pesos con geometría fija estudian esfuerzo; centros libres sin controlador estudian geometría ordinaria. No se afirma identidad exacta con algoritmos publicados completos.

Con $J_c$ candidatos cada $H$ pasos, el costo extra amortizado es aproximadamente $(C_s+J_c(C_{\rm tr}+K_pC_{\rm paso})+C_{\rm ver})/H$. Puede reutilizarse la rama de control desde el mismo estado. Se registran tiempo, memoria, parámetros y derivadas: igualar épocas no iguala trabajo. PIKAN proporciona un sustrato para observar estos cambios, sin garantizar por arquitectura su utilidad o localidad física.
\begin{figure}[htbp]\centering\includegraphics[width=\linewidth]{figuras/controlador.png}\caption{Controlador propuesto y separación entre entrenamiento, diagnóstico y prueba final. Presupuesto y conservación condicionan aceptar un cambio.}\label{fig:controlador}\end{figure}

\section{Banco de problemas y referencias}
Los casos se ordenan por estructura distribuida, localizada fija, móvil, ligada a geometría y entre campos. Cada instancia declara ecuación, condiciones, referencia y propósito. Las soluciones exactas siguientes son construcciones para verificación, no resultados entrenados.

\subsection{Poisson distribuido y localizado}
En $\Omega=(0,1)^2$, $-\Delta u=f$, $u|_{\partial\Omega}=0$:
\begin{equation}u_*=\sin(k\pi x)\sin(k\pi y),\qquad f=2(k\pi)^2u_*,\qquad k\in\{1,4,8\}.\end{equation}
Derivar dos veces cada seno obtiene $f$. $k=1$ verifica; $k=4,8$ exigen resolución distribuida. Las derivadas no son constantes, por lo que mejorar aquí no refuta automáticamente focalización. Se controla ausencia de región privilegiada persistente junto a señales permutadas y costo comparable.

En $(0,1)$, un caso localizado usa $u_*=q e^{-ad^2}$, $q=x(1-x)$, $d=x-c$, $a=100$:
\begin{equation}f=-u_*''=e^{-ad^2}\{2+4ad(1-2x)+(2a-4a^2d^2)q\}.\end{equation}
Se selecciona con $c=0.5$ y se reservan $c=0.25,0.75$ para comprobar que sigue estructura y no posición. Las fronteras son cero.

\subsection{Burgers: frente estacionario}
En $[-1,1]$, $uu_x=\nu u_{xx}$ admite $u_*=-\tanh(x/(2\nu))$, con valores exactos en frontera. Si $q=\tanh(x/(2\nu))$, $u_x=-(1-q^2)/(2\nu)$ y $u_{xx}=q(1-q^2)/(2\nu^2)$, de donde el residual es cero. Se usan $\nu\in\{0.1,0.05,0.02\}$ y región crítica $|x|\leq4\nu$. Evalúa focalización sostenida y resolución subgrilla. La estacionariedad pertenece a esta instancia, no a toda Burgers \cite{raissi2019}.

\subsection{Allen--Cahn: movimiento manufacturado}
En $[-1,1]\times[0,1]$:
\begin{equation}u_t-\varepsilon^2u_{xx}+u^3-u=f,\quad u_* =\tanh\!\left(\frac{x-c(t)}{\sqrt2\varepsilon}\right),\quad c(t)=-0.5+t.\end{equation}
Para $q=u_*$, $u_x=(1-q^2)/(\sqrt2\varepsilon)$, $u_{xx}=-q(1-q^2)/\varepsilon^2$ y $u_t=-c'(t)(1-q^2)/(\sqrt2\varepsilon)$. Difusión y reacción se cancelan, por lo que
\begin{equation}f=-c'(t)(1-u_*^2)/(\sqrt2\varepsilon).\end{equation}
Se imponen trazas exactas iniciales y de frontera, $\varepsilon\in\{0.1,0.05\}$; región crítica $|x-c(t)|\leq3\varepsilon$. El forzamiento fabrica movimiento: no se atribuye al problema homogéneo. En entrenamiento global la trayectoria es fija en espacio--tiempo; para probar liberación temporal se usan ventanas consecutivas de longitud $0.1$, comunes a los métodos, con memoria y medición del error inicial propagado.
\begin{figure}[htbp]\centering\includegraphics[width=\linewidth]{figuras/soluciones_exactas.png}\caption{Referencias analíticas de oscilación distribuida y frentes fijo y móvil. Curvas calculadas directamente, sin predicciones entrenadas.}\label{fig:soluciones}\end{figure}

\subsection{Navier--Stokes: verificación y geometría}
Las ecuaciones incompresibles son
\begin{equation}\partial_t\mathbf v+(\mathbf v\cdot\nabla)\mathbf v+\nabla p-Re^{-1}\Delta\mathbf v=\mathbf f,\qquad\nabla\cdot\mathbf v=0.\end{equation}
Kovasznay verifica el sistema estacionario con $Re=40$, $\lambda=Re/2-\sqrt{Re^2/4+4\pi^2}$, en $[-0.5,1]\times[-0.5,1.5]$:
\begin{equation}v_x=1-e^{\lambda x}\cos(2\pi y),\quad v_y=\frac\lambda{2\pi}e^{\lambda x}\sin(2\pi y),\quad p=\tfrac12(1-e^{2\lambda x}).\end{equation}
$\partial_xv_x+\partial_yv_y=0$ y $\lambda^2-Re\lambda-4\pi^2=0$ verifican las ecuaciones con $\mathbf f=0$. Se imponen velocidades exactas y una referencia de presión. Verifica acoplamiento, pero no representa un obstáculo \cite{jin2021nsfnets}.

El caso geométrico será canal $(0,2.2)\times(0,0.41)$ con cilindro de radio $0.05$ en $(0.2,0.2)$; entrada $v_x=6\bar U y(0.41-y)/0.41^2$, $v_y=0$, $\bar U=0.2$; $\nu=0.001$, $Re=\bar U(0.1)/\nu=20$. Paredes y cilindro sin deslizamiento; salida de tracción nula usando $-pI+2\nu\epsilon(\mathbf v)$. Referencia FEM mixta Taylor--Hood $P_2/P_1$ con convergencia de malla. Se declara esta instancia propuesta sin atribuirla como reproducción literal de otro banco. Turbulencia y alto Reynolds quedan fuera.

\subsection{Piezoelectricidad: derivación y transmisión}
Se estudia elasticidad lineal y electrostática sin inercia. Para desplazamiento $\mathbf d$, $S=\operatorname{sym}\nabla\mathbf d$ y $\mathbf E=-\nabla\varphi$, la entalpía local
\begin{equation}H(S,\mathbf E)=\tfrac12S:C^E:S-\mathbf E\cdot(e:S)-\tfrac12\mathbf E\cdot\kappa^S\mathbf E\end{equation}
produce, al derivar en $S$ y en $-\mathbf E$, $\sigma=C^E:S-e^T\mathbf E$ y $\mathbf D=e:S+\kappa^S\mathbf E$. Los balances son
\begin{equation}\nabla\cdot\sigma+\mathbf b=0,\qquad\nabla\cdot\mathbf D=\rho_f.\label{eq:piezo}\end{equation}
En deformación plana, $S_V=(d_{x,x},d_{y,y},d_{x,y}+d_{y,x})^T$ y $g=\nabla\varphi$:
\begin{equation}
\begin{bmatrix}\sigma_V\\\mathbf D\end{bmatrix}=\begin{bmatrix}C&e^T\\e&-\kappa\end{bmatrix}\begin{bmatrix}S_V\\g\end{bmatrix},\quad
C=\begin{bmatrix}C_{11}&C_{12}&C_{16}\\C_{12}&C_{22}&C_{26}\\C_{16}&C_{26}&C_{66}\end{bmatrix},\quad
e=\begin{bmatrix}e_{11}&e_{12}&e_{16}\\e_{21}&e_{22}&e_{26}\end{bmatrix}.
\end{equation}
$\sigma_V=(\sigma_{xx},\sigma_{yy},\sigma_{xy})^T$, $\kappa$ es $2\times2$ simétrica. El tercer componente de deformación usa cortante ingenieril. Tensión plana requiere condensar el constitutivo 3D; las tablas no se intercambian sin transformación. El antecedente local de vigas bimorfas aporta formulación y datos, cuyos signos, unidades y reducción dimensional deben verificarse antes de reproducir PVDF \cite{gonzalez2026pinnpiezo,lahmer2008}.

Integrar balances con pruebas $\mathbf v,\psi$ nulas en fronteras esenciales produce
\begin{align}\int_\Omega\epsilon(\mathbf v):\sigma\,dx&=\int_\Omega\mathbf v\cdot\mathbf b\,dx+\int_{\Gamma_t}\mathbf v\cdot\bar{\mathbf t}\,ds,\\
\int_\Omega\nabla\psi\cdot\mathbf D\,dx&=\int_{\Gamma_q}\psi\bar q\,ds-\int_\Omega\psi\rho_f\,dx.
\end{align}
El bloque eléctrico negativo define un problema de punto silla, no una energía positiva para minimizar libremente en ambos campos. Sin cargas ni fuerzas superficiales, la interfaz satisface $[\mathbf d]=[\varphi]=0$, $[\sigma n]=[\mathbf D\cdot n]=0$, con normal común. Los gradientes pueden saltar. Las pérdidas fuertes se calculan dentro de cada material y la transmisión por separado; no se derivan coeficientes discontinuos como si fueran suaves.
\begin{figure}[htbp]\centering\includegraphics[width=\linewidth]{figuras/piezo_esquema.png}\caption{Dominio bimorfo y acoplamiento constitutivo. Esquema propuesto, sin resultados FEM piezoeléctricos.}\label{fig:piezo}\end{figure}

El primer caso acoplado será adimensional y reproducible: $\Omega=(0,1)\times(-0.05,0.05)$, $C=\left[\begin{smallmatrix}2&1&0\\1&2&0\\0&0&0.5\end{smallmatrix}\right]$, $\kappa=I_2$, $e^+=\left[\begin{smallmatrix}0&0&0.1\\0.2&0.1&0\end{smallmatrix}\right]$, $e^-=-e^+$. Desplazamiento cero a la izquierda, tracción $(0,1)$ a la derecha, caras horizontales libres; $\varphi=0$ abajo, $\varphi=1$ arriba, flujo eléctrico cero en extremos; $\mathbf b=0$, $\rho_f=0$. Son materiales sintéticos, no constantes PVDF. FEM conforme a interfaz verificará campos y balances.

El caso dimensional usará la geometría publicada $L=0.1\,\mathrm m$, altura $0.001\,\mathrm m$, una vez verificados tensor, polarización, espesor fuera del plano y condiciones eléctricas. Una fuerza de $1\,\mathrm N$ no fija tracción 2D sin espesor. Esto es requisito de reproducción dimensional, no un dato inferible de una imagen.

\subsection{Predicciones previas al desarrollo}
En Poisson distribuido se espera ausencia de preferencia persistente tras controlar costo; en Burgers, concentración sostenida; en Allen--Cahn por ventanas, seguimiento sin olvido excesivo; en flujo, utilidad en regiones geométricas compatible con continuidad; en piezoelectricidad, mejora del campo peor resuelto sin perjudicar transmisión u otro campo. Son predicciones falsables. La metodología fija los controles y criterios antes de desarrollar el mecanismo.
